%
% Halaman Abstract
%
% @author  Andreas Febrian
% @version 1.00
%

\chapter*{\fontsize{14pt}{16pt}\selectfont \textbf{EXTENDED ABSTRACT}}
\vspace{-0.5cm}
\begin{center}
	\begin{spacing}{1.5}
	{\fontsize{14pt}{16pt}\selectfont \textbf{Undergraduate Thesis}}\\ [0.3cm]
	% delete the irrelevant option:Undergraduate Thesis/Master Thesis/Dissertation
	{\fontsize{14pt}{16pt}\selectfont \textbf{\judulInggris}}\end{spacing}
\end{center}
\vspace{-2cm}
\begin{singlespace}
		\begin{center}
		\bo{Student Name} : \penulis \\ 
		\bo{Student ID Number (NPM)} : {\npm}\\
		\bo{Study Program} : \jenjanginggris~Study Program of~\bidanginggris
	\end{center}
\end{singlespace}




\vspace*{0cm}

Body/Text of the abstract: {\bf 1750 words and maximum 4 pages}, with a few subheadings (Font: Times New Roman, 12, justified, Line Space 1.5)

\noindent \paragraph{Paragraph 1. Research Background, Objectives, and Methodology}

\textbf{Background and Novelty:} Briefly state the research gap and the urgency of the study.

\noindent \textbf{Objectives:} State the primary goal.

\noindent \textbf{Methodology \& Schematic Diagram:} Describe the research design and data analysis briefly. To ensure the logic is understood by English speakers, the following schematic diagram illustrates the complete workflow.

\begin{figure}[h]
	\centering
	% \includegraphics[width=0.8\textwidth]{diagram.png}
	\caption{Schematic Diagram of Research Framework}
\end{figure}

\paragraph{Paragraphs 2--4. Results and Discussion (70\% of the Extended Abstract)}

Provide a detailed, multi-paragraph technical explanation. This is the primary section for readers to understand the substance of the Indonesian-language dissertation.

\paragraph{Paragraph 2. Specific Findings}

Detail the primary data obtained. Do not merely state the results; explain the technical significance of the findings.

\paragraph{Paragraph 3. Analysis and Interpretation}

Explain the underlying reasons behind the results. Compare the findings extensively with international literature. Discuss how the results confirm or challenge existing global theories.

\paragraph{Paragraph 4. Scientific Implications}

Describe the practical and theoretical implications of the findings. Use this section to demonstrate the depth of the dissertation’s contribution to the field.

A maximum of two figures and/or tables may be included in this section.

\begin{figure}[h]
	\centering
	% \includegraphics[width=0.75\textwidth]{figure2.png}
	\caption{Figure Header}
	\small Source: Author surname and publication date (e.g., Widakdo, 2022)
\end{figure}

%\begin{table}[ht]
%	\centering
%	\caption{Summary of Rating Dimensions with Means (Mean), Standard Deviation (SD), Range, and Median}
%	\small Source: Author surname and publication date (e.g., Widakdo, 2022)
	
%	\vspace{0.2cm}
	
%	\begin{tabular}{lcccc}
%		\hline
%		\textbf{Dimension} & \textbf{Mean} & \textbf{SD} & \textbf{Range} & \textbf{Median} \\
%		\hline
%		Valence          & 4.97 & 1.79 & 2.08--7.92 & 5.20 \\
%		Arousal          & 4.86 & 2.04 & 3.04--7.28 & 4.79 \\
%		\hline
%	\end{tabular}
%\end{table}

\paragraph{Paragraph 5. Conclusion}

Summarize the main conclusions of the dissertation and highlight the broader impact or future potential of the research.

\bigskip

\noindent %\lipsum[1]

\vspace*{0.5cm}

\noindent \bo{Keywords}: Interactive Visualization, K-Means Clustering, Panel Data, PCA, Three-Way Data Analysis\\


\newpage 